\chapter*{Para saber más}
\addcontentsline{toc}{chapter}{Para saber más}
\markboth{Para saber más}{Para saber más}

Lo que sigue no es una bibliografía académica, sino una selección de lecturas para quien quiera seguir tirando del hilo. He incluido solo los textos que me parecen más accesibles o más importantes, con una breve nota sobre cada uno. Las referencias completas de los estudios científicos citados en las notas a pie de página se encuentran en las propias notas.

\section*{Sobre la biología del vínculo}

\begin{description}[style=nextline, leftmargin=0pt, labelwidth=0pt]

\item[Helen Fisher, \textit{Why We Love} (2004)]
La mejor introducción a la neurociencia del enamoramiento. Fisher explica con claridad cómo la dopamina, la noradrenalina y la serotonina crean ese estado de «locura transitoria» que llamamos enamorarnos. Muy ameno y riguroso a la vez.

\item[Helen Fisher, \textit{Anatomy of Love} (1992, ed. rev. 2016)]
Más ambicioso que el anterior: recorre la historia evolutiva del emparejamiento humano, desde la sabana africana hasta las apps de citas. Imprescindible para entender por qué somos como somos en pareja.

\item[John Bowlby, \textit{Attachment and Loss} (1969--1980)]
La obra fundacional de la teoría del apego. Son tres volúmenes densos, pero el primero (\textit{Attachment}) es el más accesible y el que más directamente conecta con lo que contamos en los capítulos 1 y 2. Para quien prefiera algo más breve, el artículo de Hazan y Shaver (1987), que aplica la teoría del apego al amor de pareja, es un buen punto de partida.

\item[Richard Wrangham, \textit{Catching Fire} (2009)]
Un libro fascinante sobre cómo el dominio del fuego y la cocción de alimentos cambiaron el curso de la evolución humana. Lo citamos brevemente en el capítulo 1, pero merece una lectura completa: cambia la forma de ver nuestra propia especie.

\item[Ashley Montagu, \textit{Touching} (1971)]
Un clásico sobre el sentido del tacto y su papel en el desarrollo humano. Sorprende descubrir cuánto depende nuestra salud física y emocional del contacto piel con piel.

\end{description}

\section*{Sobre el amor y la pareja}

\begin{description}[style=nextline, leftmargin=0pt, labelwidth=0pt]

\item[Erich Fromm, \textit{El arte de amar} (1956)]
Un pequeño gran libro. Fromm argumenta que el amor no es un sentimiento que te sucede, sino un arte que se aprende y se practica. Su distinción entre amor inmaduro («te amo porque te necesito») y amor maduro («te necesito porque te amo») sigue siendo una de las mejores frases escritas sobre el tema.

\item[C.\,S. Lewis, \textit{Los cuatro amores} (1960)]
Lewis distingue entre afecto, amistad, eros y caridad con una lucidez y un humor que solo él tenía. Es breve, profundo y muy disfrutable. Ideal para leer en pareja.

\item[Deborah Tannen, \textit{You Just Don't Understand} (1990)]
Un estudio sobre cómo hombres y mujeres usan el lenguaje de formas distintas, y cómo eso genera malentendidos crónicos en las parejas. Práctico y revelador. Existe traducción al español.

\end{description}

\section*{Sobre la visión cristiana del matrimonio}

\begin{description}[style=nextline, leftmargin=0pt, labelwidth=0pt]

\item[Juan Pablo II, \textit{Hombre y mujer los creó} (1979--1984)]
Las catequesis de los miércoles sobre la «Teología del Cuerpo». Es la obra más ambiciosa jamás escrita sobre el significado del cuerpo, la sexualidad y el matrimonio desde una perspectiva cristiana. No es lectura fácil ---el estilo de Wojtyła es filosófico y denso---, pero las ideas que contiene son revolucionarias. Para una introducción más accesible, véase el libro de Yves Semen que se cita más abajo.

\item[Yves Semen, \textit{La sexualidad según Juan Pablo II} (2005)]
La mejor puerta de entrada a la Teología del Cuerpo. Semen traduce el pensamiento de Wojtyła a un lenguaje asequible sin traicionar su profundidad. Muy recomendable como primer contacto.

\item[Benedicto XVI, \textit{Deus Caritas Est} (2005)]
La primera encíclica de Benedicto XVI es una meditación luminosa sobre el amor: desde el eros hasta el ágape, desde la experiencia humana hasta la caridad de Dios. Breve, bella y sorprendentemente personal.

\item[Francisco, \textit{Amoris Laetitia} (2016)]
La exhortación del Papa Francisco sobre la familia. El capítulo 4, que es un comentario al himno de la caridad de san Pablo (1~Cor 13), es una de las páginas más hermosas escritas sobre el amor conyugal en un documento eclesial. Muy pastoral, muy realista, muy humana.

\item[Juan Pablo II, \textit{Familiaris Consortio} (1981)]
El documento de referencia sobre la familia en la doctrina católica contemporánea. Más sistemático que \textit{Amoris Laetitia}, pero complementario. Especialmente valioso para entender la indisolubilidad, la fidelidad y la apertura a la vida.

\item[\textit{Gaudium et Spes}, nn. 47--52 (1965)]
Los números del Concilio Vaticano II dedicados al matrimonio y la familia siguen siendo un texto de referencia. Breve, denso y profundamente renovador para su época.

\end{description}

\section*{Para pensar más}

\begin{description}[style=nextline, leftmargin=0pt, labelwidth=0pt]

\item[Emmanuel Levinas, \textit{Totalidad e infinito} (1961)]
Para lectores filosóficamente audaces. Levinas desarrolla la idea de que el encuentro con el «rostro del otro» es la experiencia ética fundamental. Su concepto de alteridad ---el otro como irreductiblemente otro--- ilumina de forma profunda lo que decimos en el capítulo 5 sobre la diferencia como condición del encuentro.

\item[Gabriel Marcel, \textit{Être et avoir} (1935)]
Marcel distingue entre «tener» y «ser», entre poseer y entregarse. Su filosofía del compromiso y la fidelidad creadora conecta directamente con lo que exploramos en los capítulos 8 y 9. Breve y luminoso.

\end{description}
